\documentclass[11pt]{article}
\usepackage[margin=1in]{geometry}
\usepackage{amsmath,amssymb,amsthm}
\usepackage{graphicx}
\usepackage[hidelinks]{hyperref}

\newtheorem{theorem}{Theorem}
\newtheorem{corollary}[theorem]{Corollary}

\newcommand{\PP}{\mathcal P}
\newcommand{\BB}{\mathcal B}
\newcommand{\Nat}{\mathbb N}
\newcommand{\Rat}{\mathbb Q}
\newcommand{\Span}{\operatorname{span}}

\begin{document}

\begin{center}
{\Large A fixed-representative partial reduction for Question 1 of arXiv:1912.03994}

\medskip
{\small Run \texttt{fa\_banach\_001}; source arXiv:1912.03994}
\end{center}

\section*{Status}

\textbf{partial result.}

Cuth--Dolezal--Doucha--Kurka ask whether there is a continuous map
\(\Phi:\PP_\infty\to\BB\) such that \(X_\mu\) is linearly isometric to
\(X_{\Phi(\mu)}\) for every \(\mu\in\PP_\infty\). This packet gives a
positive answer on every stratum where an independent dense representative
sequence is fixed in advance.

\section*{Source Questions}

The relevant source questions appear in arXiv:1912.03994, pages 20--21.

\begin{center}
\includegraphics[width=0.95\linewidth]{figures/open_problem_crop_page20.png}

\medskip
\includegraphics[width=0.95\linewidth]{figures/open_problem_crop_page21.png}
\end{center}

Question 1 asks for a continuous reduction \(\PP_\infty\to\BB\). Question 2
asks for the corresponding \(\boldsymbol{\Sigma}^0_2\)-measurable reduction
from \(\PP_\infty\) to \(SB(X)\) with an admissible topology.

\section*{Definitions}

Let \(V\) be the rational vector space with basis \((e_n)_{n\in\Nat}\), as in
the source paper. For \(\mu\in\PP\), write \(X_\mu\) for the completion of
\((c_{00},\mu)\) modulo its null space, and write \([v]_\mu\) for the image of
\(v\in V\) in \(X_\mu\).

Fix a strictly increasing sequence \(s=(s_j)_{j\ge1}\) in \(\Nat\). Let
\(\PP_s\subseteq\PP_\infty\) be the class of all \(\mu\) such that:
\begin{enumerate}
\item for every nonzero finite rational scalar family \((a_j)\),
\[
\mu\Big(\sum_j a_j e_{s_j}\Big)>0,
\]
so the vectors \([e_{s_j}]_\mu\) are linearly independent;
\item the closed linear span of \(\{[e_{s_j}]_\mu:j\ge1\}\) is \(X_\mu\).
\end{enumerate}

\section*{Proof Intuition}

The source paper already gives a \(\boldsymbol{\Sigma}^0_2\)-measurable
reduction by choosing, for each \(\mu\), the first coordinates that remain
linearly independent in \(X_\mu\). The discontinuity can only enter through
that coordinate-selection step. If the representative coordinates are fixed
beforehand, no selection remains: the output norm is just the restriction of
\(\mu\) to that fixed subsequence. Continuity is then immediate from the
coordinate topology on \(\PP\).

\section*{Theorem}

\begin{theorem}
For every strictly increasing sequence \(s=(s_j)\) in \(\Nat\), the formula
\[
\Phi_s(\mu)\Big(\sum_{j=1}^m a_j e_j\Big)
 =
\mu\Big(\sum_{j=1}^m a_j e_{s_j}\Big),
\qquad a_j\in\Rat,
\]
defines a continuous map \(\Phi_s:\PP_s\to\BB\), where \(\PP_s\) has the
relative topology inherited from \(\PP_\infty\). Moreover
\[
X_{\Phi_s(\mu)}\equiv X_\mu
\]
linearly isometrically for every \(\mu\in\PP_s\).
\end{theorem}

\begin{proof}
Fix \(s\) and \(\mu\in\PP_s\). The formula defines a pseudonorm on \(V\)
because it is the pullback of the pseudonorm \(\mu\) under the fixed rational
linear injection \(T_s(e_j)=e_{s_j}\). The first defining property of
\(\PP_s\) says exactly that no nonzero vector of \(V\) is sent to a
\(\mu\)-null vector; hence \(\Phi_s(\mu)\) is a norm and belongs to \(\BB\).

For continuity, let \(v=\sum_{j=1}^m a_j e_j\in V\). Then
\[
\Phi_s(\mu)(v)=\mu(T_s v).
\]
Since \(T_s v\in V\) is fixed, the map \(\mu\mapsto \mu(T_s v)\) is one of
the coordinate maps defining the topology of \(\PP\). Therefore every
coordinate of \(\Phi_s\) is continuous, which is precisely continuity as a map
into \(\BB\subseteq\PP\).

Finally define \(J_\mu:V\to X_\mu\) by \(J_\mu(e_j)=[e_{s_j}]_\mu\). For
every \(v\in V\),
\[
\|J_\mu v\|_{X_\mu}=\Phi_s(\mu)(v).
\]
Thus \(J_\mu\) descends to an isometric embedding of \(X_{\Phi_s(\mu)}\) into
\(X_\mu\). Its range contains the span of \(\{[e_{s_j}]_\mu:j\ge1\}\), which
is dense in \(X_\mu\) by the second defining property of \(\PP_s\). Hence the
extension of \(J_\mu\) is onto and is a linear isometry.
\end{proof}

\begin{corollary}
Let \(A\subseteq\Nat\) have infinite complement \(S=\{s_1<s_2<\cdots\}\). On
the stratum of all \(\mu\in\PP_\infty\) for which the coordinates
\(\{e_i:i\in A\}\) vanish and the quotient vectors
\(\{[e_{s_j}]_\mu:j\ge1\}\) are linearly independent, Question 1 has a
positive answer via \(\Phi_s\).
\end{corollary}

\begin{proof}
Under the stated hypotheses, every coordinate outside \(S\) is zero in the
quotient and every coordinate inside \(S\) belongs to the fixed representative
sequence. Hence the span of \([e_{s_j}]_\mu\) is dense in \(X_\mu\), and the
theorem applies.
\end{proof}

\section*{Relation to Question 2}

The source paper proves a \(\boldsymbol{\Sigma}^0_2\)-measurable reduction
from \(\BB\) to \(SB_\infty(X)\) for any admissible topology on \(SB(X)\).
Composing that reduction with the continuous map \(\Phi_s\) gives the
corresponding relative positive answer to Question 2 on each stratum
\(\PP_s\).

\section*{Verification and Limitations}

The proof is purely formal and uses only the definitions of the source coding
spaces. It does not solve the global problem because it assumes that the
representative sequence \(s\) is fixed on the whole domain. Thus the remaining
obstruction is precisely the source paper's global selection problem: choosing
such an independent dense sequence continuously as \(\mu\) varies over all of
\(\PP_\infty\).

Local index searches for arXiv:1912.03994, Question 1, \(\PP_\infty\),
\(\BB\), continuous reduction, and admissible topologies found no prior packet
for this fixed-representative partial result. Web searches on 2026-06-16 found
the source paper and the sequel arXiv:2204.06834, but no exact global answer
to Question 1 or Question 2.

\section*{References}

\begin{thebibliography}{9}

\bibitem{CDDK}
M. Cuth, M. Dolezal, M. Doucha, and O. Kurka,
\emph{Polish spaces of Banach spaces},
arXiv:1912.03994.

\bibitem{CDDK2}
M. Cuth, M. Dolezal, M. Doucha, and O. Kurka,
\emph{Polish spaces of Banach spaces. Complexity of isometry and isomorphism classes},
arXiv:2204.06834.

\end{thebibliography}

\end{document}

